\documentclass[11pt]{article}

% Cambiar "review" por "final" para la versión final.
\usepackage[review]{acl}

\usepackage{times}
\usepackage{latexsym}
\usepackage[T1]{fontenc}
\usepackage[utf8]{inputenc}
\usepackage[spanish,es-noshorthands]{babel}
\usepackage{microtype}
\usepackage{inconsolata}
\usepackage{graphicx}
\usepackage{amsmath}
\usepackage{booktabs}

\title{Ontogenia Artificial: trayectorias de adquisición sintáctica en Pythia y su alineación con la adquisición humana del lenguaje}

\author{
  Sebastián Mesch Henriques \\
  UDESA \\
  \texttt{sebastianmeschhenriques@gmail.com} \\\And
  Leandro \\
  UDESA \\
  \texttt{\_} \\\And
  Hugo \\
  UDESA \\
  \texttt{\_} \\
}

\begin{document}
\maketitle

\begin{abstract}
\textit{Placeholder del abstract (150--200 palabras). Se completa en Semana 4.}

Evaluamos longitudinalmente los checkpoints intermedios de la suite Pythia sobre pares mínimos sintácticos (BLiMP y Zorro) para caracterizar la topología de las trayectorias de aprendizaje a nivel de paradigma individual, y correlacionamos el orden de estabilización con datos de adquisición humana (Wordbank). Mostramos que (i) una fracción no trivial de los paradigmas exhibe curvas no monótonas, (ii) el ranking de emergencia correlaciona con la edad de adquisición humana, y (iii) los paradigmas con excepciones morfológicas presentan curvas en U más profundas que los puramente estructurales.
\end{abstract}

\section{Introducción}
\label{sec:introduccion}

\textit{Placeholder. Motivación, gap, contribución, roadmap del paper.}

\section{Trabajo Relacionado}
\label{sec:related}

\textit{Developmental interpretability, curvas en U sintácticas, alineación humano-modelo.}

\section{Método}
\label{sec:metodo}

\subsection{Modelos y checkpoints}
\textit{Pythia 14m/160m/410m deduplicados, 24 checkpoints log-espaciados.}

\subsection{Tareas y métrica}
\textit{BLiMP, Zorro. SLLN-LP como métrica principal.}

\subsection{Alineación humana}
\textit{Wordbank como fuente primaria de AoA; literatura complementaria para sintaxis abstracta.}

\section{Experimentos}
\label{sec:experimentos}

\textit{Grilla experimental, hipótesis operacionalizadas H1--H3.}

\section{Resultados}
\label{sec:resultados}

\textit{Figuras de trayectorias, clasificación topológica, correlación Spearman con bootstrap.}

\section{Discusión}
\label{sec:discusion}

\textit{Interpretación, limitaciones, líneas futuras (SLT/LLC, extensión a léxico/EWoK).}

\section{Conclusión}
\label{sec:conclusion}

\textit{Placeholder.}

\section*{Limitaciones}

\textit{Sección requerida por política ACL desde 2023.}

\bibliography{custom}
\bibliographystyle{acl_natbib}

\appendix

\section{Detalles adicionales}
\label{sec:apendice}

\textit{Tablas completas, hiperparámetros, curvas por paradigma.}

\end{document}
